\documentclass[../../main2.tex]{subfiles}

\begin{document}

	% ======================================================================
	% 骨架文件：小节标题与追问链为终版；正文由后续会话填充。
	% 硬要求（notes/arena-prompt.md 任务卡 D）：
	%   §13.3 必须真实推导一个具体例子（不许空话）；
	%   §13.4 失败模式清单必须列满 5 条以上。
	% ======================================================================

	\chapter{从贡献表到未来走向}
	\label{ch:forecast}

	\begin{motivation}
		上一章遗留的问题：四条通道全部定义完毕，贡献分解的公式在手——但那是对\textbf{过去}的分解。我们手里拿着 $C_{\text{模因}}, C_{\text{经济}}, C_{\text{政治}}, C_{\text{法律}}$ 四个数字，怎么用它说出「接下来会怎样」？\\
		\textbf{核心动机}：把贡献表拼回时间线——重组条件概率链，给出一条可辩护的未来。这是全书唯一一次完整演示「如何用这本书」。
	\end{motivation}

	\begin{logicmap}
	\begin{tikzpicture}[node distance=0.9cm, every node/.style={draw, rectangle, rounded corners, align=center}]
	\node (q) {贡献表怎么变成「接下来」？};
	\node (s) [below=of q] {贡献表 = 状态压缩（§13.1）};
	\node (a) [below=of s] {条件概率链重组（§13.2）};
	\node (e) [below=of a] {一个完整演示（§13.3）};
	\node (f) [below=of e] {失败模式清单（§13.4）};
	\draw[-{Stealth}] (q) -- (s);
	\draw[-{Stealth}] (s) -- (a);
	\draw[-{Stealth}] (a) -- (e);
	\draw[-{Stealth}] (e) -- (f);
	\end{tikzpicture}
	\end{logicmap}

	% ==================================================================
	\section{贡献表是一份状态压缩}
	\label{sec:fc-table}

	\textcolor{gray}{\textit{【骨架 · 追问链（终版）】把四个通道的贡献看成状态变量。诚实说明：这是粗粒度压缩，会丢信息。丢了什么？→ 丢个体差异、丢通道内部的异质性——一张表描述不了一亿人。}}

	\textcolor{gray}{（正文待写：压缩的取舍论证（呼应第2章地图类比）；「能接受多少丢失」取决于预测目标的时间尺度与粒度。）}

	\begin{readingnote}
		\textcolor{gray}{（待写：你有没有想过，为什么天气预报不说「每一滴雨落在哪里」？）}
	\end{readingnote}

	\paragraph*{逻辑衔接}
	\textcolor{gray}{（待写）压缩后的状态要能重新展开成未来。下一节「条件概率链的重新组装」把第6章的因子换成通道贡献。}

	% ==================================================================
	\section{条件概率链的重新组装}
	\label{sec:fc-reassemble}

	\textcolor{gray}{\textit{【骨架 · 追问链（终版）】用第6章的乘积分解，但把因子换成通道贡献。这不就是第6章吗？→ 不是：第6章是概念，这里是可执行——每个因子 $P(C_k \mid C_{k-1})$ 都有对应的通道贡献与测量接口。}}

	\textcolor{gray}{（正文待写：组装的完整流程图；与附录 A.2 的接口；「可执行」与「概念」的差别：数据来源、更新规则、误差传播。）}

	\paragraph*{逻辑衔接}
	\textcolor{gray}{（待写）流程有了，但它必须在一个真实事件上跑通一次才算数。下一节「一个具体工作过的例子」做全书唯一一次完整演示。}

	% ==================================================================
	\section{一个具体工作过的例子}
	\label{sec:fc-example}

	\textcolor{gray}{\textit{【骨架 · 追问链（终版）】必须真实推导一个例子（如：某项政策出台后，某社会指标的概率分布）——逐步算四通道的贡献分解，展示怎么得到概率判断。为什么需要例子？→ 因为抽象方法无法自证。}}

	\textcolor{gray}{（正文待写（后续会话硬要求）：选一个读者熟悉的社会事件；逐步代入：状态刻画 $\to$ 四通道 $C$ 的估计 $\to$ 链式组装 $\to$ 概率断言 + 依赖因素清单；诚实标注每一步的假设与数据来源。\lowfreq{例子的具体数字是示意性估计，不是实证结果。}）}

	\begin{questionbox}
		\textcolor{gray}{（待写：反驳位——「一个例子证明不了方法有效」。回答：对——例子演示的是「可执行性」不是「有效性」；有效性要看失败模式清单（§13.4）与方法边界。）}
	\end{questionbox}

	\paragraph*{逻辑衔接}
	\textcolor{gray}{（待写）演示跑通了。现在必须做另一半诚实的事：这个方法会在哪里翻车？下一节「预测失败模式清单」。}

	% ==================================================================
	\section{预测失败模式清单}
	\label{sec:fc-failure}

	\textcolor{gray}{\textit{【骨架 · 追问链（终版）】列出本书方法会失败的情况（至少 5 条）：权重不可测、交互项过大、通道漏项、外部冲击、非平稳。诚实列出失败模式——这是可信度的来源。}}

	\textcolor{gray}{（正文待写：五种失败模式逐一给出（i）症状（ii）为什么无法修复（iii）部分缓解手段；「少于 3 条的失败清单不值得信任」（引言的承诺在此兑现）。）}

	\begin{reader_anticipation}
		\item \textcolor{gray}{（待写）失败之后怎么办？→ 第14章：预测的反身性——预测本身改变被预测的世界。}
	\end{reader_anticipation}

	\paragraph*{逻辑衔接}
	\textcolor{gray}{（待写）我们会算未来了，也知道了哪里算不动。还剩最后一问：我们预测出的未来，会不会反过来改变人的动机？第14章「循环的闭合」处理这个自指。}

	% ==================================================================
	\section{本章小结与衔接}
	\label{sec:fc-summary}

	\textcolor{gray}{（正文待写：收束——从贡献表到概率断言的完整链路 + 边界。）}

	\paragraph*{逻辑衔接}
	\textcolor{gray}{（待写）第14章「循环的闭合：未来塑造新的动机」——全书的最后一章，回到第1章的起点。}

\end{document}
